% ============================================================================
% CHAPTER 5: CONCLUSION AND FUTURE SCOPE
% ============================================================================

\chapter{CONCLUSION AND FUTURE SCOPE}

\section{CONCLUSION}

This project successfully developed an advanced deep learning system for objective depression detection from EEG signals, addressing critical methodological flaws that plague existing research in this domain. The resulting system achieves clinically meaningful accuracy while providing interpretable predictions validated against established neurophysiological biomarkers of depression.

\subsection{Summary of Achievements}

The most significant achievement of this work is the design and implementation of a novel multi-branch deep learning architecture that effectively captures the complexity of depression-related EEG patterns. The hybrid model combines a Vision Transformer for analyzing time-frequency representations with a Graph Attention Network for modeling spatial electrode relationships, connected through an attention-based fusion mechanism that adaptively weights their complementary contributions. This architecture represents a substantial advance over the simple convolutional networks used in prior work, enabling the model to capture both global temporal dynamics through Transformer self-attention and local spatial connectivity patterns through graph message passing.

The second major achievement is the implementation of an honest evaluation protocol that eliminates the data leakage problem endemic to prior research in this field. By employing Leave-One-Subject-Out cross-validation with complete subject-level separation between training and test sets, the reported 91.38\% subject-level accuracy represents true generalization capability to completely unseen individuals. This accuracy, while numerically lower than the inflated metrics reported by studies with leaky validation, reflects the actual performance that would be observed in clinical deployment and is therefore far more valuable for translation to practice.

The third achievement is comprehensive clinical validation through multi-level explainability analysis. The combination of Integrated Gradients for feature attribution, Layer-wise Relevance Propagation for understanding information flow, and Testing with Concept Activation Vectors for validating high-level clinical concepts provides unprecedented insight into model decision-making. This analysis confirmed that the model learns established depression biomarkers including frontal alpha asymmetry, elevated frontal theta activity, and delta abnormalities, rather than artifacts or subject-specific patterns that would not generalize.

Finally, the memory-efficient implementation makes the system accessible for broader research use without requiring specialized computing infrastructure. Through mixed-precision training, gradient accumulation, and careful architecture design, the system trains successfully on consumer-grade GPUs with 8GB VRAM, democratizing access to advanced deep learning approaches for EEG-based psychiatric research.

\subsection{Key Contributions}

This work makes substantial contributions across technical, methodological, and clinical dimensions. On the technical front, this project represents the first implementation combining Vision Transformer and Graph Attention Network architectures for EEG-based depression detection. The multi-wavelet Wavelet Packet Decomposition approach using Daubechies-4, Symlet-5, and Coiflet-3 wavelets provides a richer spectral representation than single-wavelet approaches, capturing diverse signal characteristics through the complementary properties of these wavelet families. The attention-based fusion mechanism enables adaptive combination of information from different processing branches, learning to emphasize whichever modality is more informative for each particular input.

The methodological contributions address fundamental problems with evaluation practices in the field. The rigorous LOSO validation protocol eliminates data leakage by ensuring complete subject-level separation, providing honest performance estimates that reflect true clinical utility. The comprehensive comparison framework developed in this work clearly demonstrates the dramatic difference between leaky and proper evaluation, highlighting why the apparent 98\% accuracy reported with standard cross-validation does not translate to real-world performance. The multi-level explainability framework combining three complementary methods provides a template for validating that models learn clinically meaningful patterns.

The clinical contributions demonstrate the potential for translating this research toward practical application. By showing that proper validation yields clinically meaningful accuracy around 91\% rather than the artificial inflation to 98\% seen with leaky methods, this work provides realistic expectations for what EEG-based depression detection can achieve. The validation of model decision-making against established EEG biomarkers builds confidence that the learned representations have genuine neurophysiological significance. Together, these contributions establish a framework for future clinical deployment with interpretable, validated predictions.

\subsection{Limitations}

While this work achieves significant improvements over prior approaches, several limitations should be acknowledged to properly contextualize the results and guide future research directions.

The dataset size of 58 subjects, while sufficient for demonstrating the approach, is relatively small by deep learning standards. Deep neural networks with over 1.5 million parameters typically benefit from larger training sets, and the generalization claims would be strengthened by validation on larger multi-site datasets encompassing hundreds or thousands of subjects. The current results may not fully capture the variability in depression presentation and EEG characteristics across broader populations.

The use of only eyes-closed resting-state recordings represents a constraint on the information available to the model. Task-based recordings during cognitive or emotional challenges, or eyes-open resting-state conditions, might provide additional discriminative information that could improve classification accuracy. Different recording conditions evoke different brain states and may reveal complementary aspects of depression-related neural dysfunction.

The binary classification framework distinguishing only Major Depressive Disorder from healthy controls, while clinically relevant for screening, provides limited nuance. Depression exists on a spectrum of severity, and different subtypes may have distinct neurobiological profiles. Extending to severity prediction or subtype classification would provide more clinically actionable information for treatment planning.

The demographic composition of the dataset may not represent all populations that would be encountered in clinical deployment. Depression prevalence, expression, and EEG characteristics may vary across age groups, ethnic backgrounds, and cultural contexts. Validation across diverse demographics is essential before clinical deployment to ensure the system works equitably for all potential patients.

Finally, the unknown or uncontrolled medication status of participants introduces potential confounds. Psychotropic medications can substantially alter EEG patterns, and the current analysis cannot distinguish medication effects from depression effects. Studies with careful medication control, either through medication-naive populations or systematic investigation of medication influences, would strengthen confidence in the findings.

\subsection{Comparison with Objectives}

Table \ref{tab:objective_achievement} summarizes the achievement of stated objectives, demonstrating that all primary targets were met or exceeded.

\begin{table}[H]
    \centering
    \small
    \caption{Achievement of Project Objectives}
    \label{tab:objective_achievement}
    \begin{tabularx}{\textwidth}{|X|L{3cm}|L{3.5cm}|}
        \hline
        \textbf{Objective} & \textbf{Target} & \textbf{Achieved} \\
        \hline
        Subject-level accuracy (LOSO) & $>$ 85\% & 91.38\% \\
        \hline
        AUC-ROC & $>$ 0.85 & 0.9042 \\
        \hline
        GPU memory requirement & $\leq$ 8 GB & $\sim$ 6 GB \\
        \hline
        Multi-level explainability & IG, LRP, TCAV & All implemented \\
        \hline
        Biomarker validation & Clinical concepts & Confirmed \\
        \hline
        No data leakage & LOSO validation & Verified \\
        \hline
    \end{tabularx}
\end{table}

\section{FUTURE SCOPE}

Several promising directions exist for extending this work across short-term, medium-term, and long-term horizons.

\subsection{Short-Term Improvements}

The immediate priority is exploring architectural enhancements that could further improve classification accuracy. A V2 model architecture adding a third branch with a Bidirectional LSTM encoder for raw EEG temporal dynamics is ready for training and evaluation. Early experiments suggest potential improvement of 2-4 percentage points in accuracy by capturing sequential patterns that the Transformer and GNN branches may miss, as the LSTM's recurrent structure is particularly suited for modeling the continuous evolution of brain states over time.

Data augmentation strategies offer another avenue for short-term improvement. Implementing time-shifting that randomly offsets epochs within the recording, amplitude scaling that varies signal magnitude across physiologically plausible ranges, and channel dropout that randomly masks individual electrodes during training could improve model robustness and reduce overfitting on the limited dataset. These augmentations simulate natural variation in recording conditions and encourage the model to learn features that generalize across such variations.

Ensemble methods combining multiple model variants through prediction averaging or voting represent a straightforward approach to improving accuracy and reducing variance. Ensembles of the current V1 model trained with different random seeds, or combinations of V1 and V2 architectures with different configurations, could provide more stable predictions than any single model while potentially capturing complementary patterns.

\subsection{Medium-Term Goals}

Multi-site validation represents a critical medium-term goal for establishing the generalizability of the approach. Validation on additional datasets from different acquisition sites using different EEG equipment and different populations would strengthen confidence that the learned patterns are truly universal markers of depression rather than artifacts of specific recording conditions. Candidate datasets include the MODMA dataset, the Child Mind Brain-Phenotype repository, and various private clinical collections that may become available through research collaborations.

Extension from binary classification to depression severity prediction would provide more clinically nuanced information than the current MDD versus healthy distinction. Regressing continuous severity scores such as PHQ-9 or BDI-II values would enable tracking of symptom trajectory and provide finer-grained information for treatment decisions. This extension requires datasets with quantitative symptom assessments rather than only diagnostic labels.

Longitudinal analysis processing multiple recordings per patient over time could enable treatment response monitoring, one of the most valuable potential clinical applications. By tracking how EEG biomarkers change during the course of therapy, clinicians could receive objective feedback on treatment effectiveness before subjective symptom improvement becomes apparent. This capability could accelerate treatment optimization and reduce the duration of inadequately treated depression.

Subtype classification distinguishing between depression subtypes such as melancholic, atypical, or anxious presentations based on EEG patterns could support personalized treatment selection. Different depression subtypes respond differently to various treatments, and biomarker-based subtyping could guide clinicians toward interventions most likely to be effective for each individual patient.

\subsection{Long-Term Vision}

The long-term vision encompasses clinical translation, regulatory approval, and expansion to broader applications. A prospective clinical trial using the system as a screening tool alongside standard clinical assessment would validate real-world clinical utility in the intended deployment context. Such a trial would recruit patients presenting with possible depression, compare model predictions with expert clinician diagnosis, and assess the added value of EEG-based biomarker information for diagnostic accuracy and confidence.

Pursuing regulatory clearance through the FDA 510(k) pathway or CE marking for the European market would enable clinical deployment as a validated medical device for depression screening. This regulatory process requires substantial clinical validation data, quality management systems, and demonstration of safety and effectiveness, representing a multi-year effort but opening the path to widespread clinical adoption.

Adaptation of the system for consumer-grade EEG headsets such as Muse or Emotiv devices could enable widespread screening outside traditional clinical settings. While such devices provide lower signal quality than research-grade systems, the model architecture could potentially be adapted or retrained for the specific characteristics of consumer hardware, enabling depression screening in primary care offices, workplaces, or even homes.

Optimization for real-time processing would enable immediate feedback during clinical consultations rather than requiring offline analysis. Real-time deployment requires careful engineering to minimize latency while maintaining accuracy, potentially through model compression techniques such as pruning, quantization, or knowledge distillation.

Extension to multi-disorder detection would create a comprehensive neural biomarker platform applicable across psychiatric conditions with EEG signatures. Anxiety disorders, attention-deficit/hyperactivity disorder, schizophrenia, and bipolar disorder all show characteristic EEG patterns that could potentially be detected using adapted versions of the current architecture, enabling a single system to screen for multiple conditions.

\subsection{Research Directions}

Beyond direct extensions of the current work, several broader research directions merit exploration. Self-supervised pre-training on large unlabeled EEG datasets using contrastive learning techniques could improve feature representations before fine-tuning on labeled depression data. The abundance of unlabeled EEG data from various research studies could provide the scale needed to learn robust general-purpose EEG features that transfer well to downstream classification tasks.

The development of large-scale EEG foundation models analogous to GPT for language or CLIP for vision represents an emerging research frontier. Such models, trained on massive and diverse EEG datasets, could enable few-shot or zero-shot learning for various neurological and psychiatric conditions, dramatically reducing the labeled data requirements for clinical applications.

Using the trained model as a tool for causal discovery could advance neuroscientific understanding of depression beyond the immediate clinical application. The attention weights and attribution maps reveal which brain regions and connections the model relies upon for classification, and systematic investigation of these patterns could generate hypotheses about causal relationships between neural circuit dysfunction and depressive symptoms.

\subsection{Conclusion}

This project demonstrates that rigorous methodology can produce EEG-based depression detection systems that are both accurate and clinically interpretable. The combination of honest evaluation through LOSO validation, advanced architecture design with complementary Transformer and GNN branches, and comprehensive explainability analysis confirming biomarker learning establishes a framework for reliable psychiatric diagnostics.

The 91.38\% subject-level accuracy achieved with LOSO represents genuine generalization capability that would be observed when deploying the model on new patients. While this figure appears numerically lower than the 98\% claimed by studies with leaky validation, it reflects the reality of clinical deployment rather than artifacts of inappropriate evaluation methodology. This distinction is crucial for translating research into clinical practice, where honest performance expectations enable appropriate clinical integration.

The framework established here---multi-branch deep learning with proper validation and multi-level explainability---provides a foundation for future work toward reliable, interpretable, and clinically deployable EEG-based psychiatric diagnostics. As larger datasets become available, validation expands to diverse populations, and regulatory pathways are navigated, systems built on these principles may contribute meaningfully to improving depression diagnosis and treatment worldwide.

